{ \setlength{\parskip}{0.7em}

\chapter[結論]{結論}\label{chapt6}

火山大国である日本において, 
火山ガスの中でも特に毒性の強い二酸化硫黄濃度の監視が重要である. 

本研究は, ラマンDIALの分布観測, 遠隔, 時間連続性に注目し, 
これによる\ce{SO2}観測を目標として, 
シミュレーションによる測定系の評価を行った. 

ラマンDIALにおける従来法では, 
ストークス散乱波長を用いて2つの受信波長を構成する
測定系が検討されていたが, 
本研究では受信波長構成を見直し, 
アンチストークス散乱波長を利用する構成と従来法を比較した. 

結果として, 
従来法に対して, 
微弱な信号強度が課題で光学系での利用が十分検討されてこなかった
アンチストークス散乱波長を含めた受信波長構成の優位性が明らかになった. 
また, ラマンDIALによる\ce{SO2}測定の検討として, 
統計誤差の抑制, \ce{H2S}干渉による影響の補正をする信号処理の検証と, 
プルームモデルに従う現実的な火山ガス分布に対する測定シミュレーション, 
干渉フィルターの設計を行った. 
それぞれ, 誤差抑制のための信号処理の定量的な評価を行い, 
局所的な高濃度分布に対する距離分解能の要求性能, 
隣接するラマン散乱スペクトルを選択するための
干渉フィルターの仕様が明らかとなった. 

今後の展望として, 
アンチストークス散乱のノイズ耐性に関する検討を挙げる. 
本研究では, アンチストークス散乱はストークス散乱に対して
信号強度が1/10と仮定し, 
諸ノイズに関するパラメータを理想形で設定した. 
その上で, 火山ガスの吸収断面積スペクトル影響下での
アンチストークス散乱を含む受信波長構成の
波長的優位性を示したが, 
一方で, 信号強度的な優位性は不明である. 
したがって, ノイズを付加したライダー信号によるラマンDIALの
\ce{SO2}濃度分布測定を改めてシミュレーションすることで
測定性能を評価する必要がある. 

% ラマンDIALの統計誤差の
% レーザ波長と受信波長構成に対する依存性に着目し, 
% 最適な波長仕様の検討を行った. 
% 結果としてレーザ波長\qty{334.6}{nm}を用いて, 
% on波長を\ce{O2}アンチストークス散乱波長の\qty{318.0}{nm} 
% off波長を\ce{O2}ストークス散乱波長の\qty{353.0}{nm}とする場合が
% 最も二酸化硫黄測定に最適なラマンDIALの仕様であり, 
% 信号強度が微弱なアンチストークス散乱も波長選択によって十分に選択可能であることが明らかとなった. 


% 距離分解能の調整による統計誤差の抑制について検証し, 
% \ce{SO2}が\qty{100}{ppm}程度の高濃度領域で要求する測定誤差を達成することがわかった.
% 干渉ガス濃度の推定による補正の効果についても検証し, 
% 霧島火山のような\ce{H2S}が極めて高濃度な領域で測定する場合, 
% その影響を完全に補正できない可能性があり, 
% 事前に妥当な\ce{H2S}濃度の推定値を用意するか, 
% 補助的に\ce{H2S}を測定する手法を併用する必要が示された. 


% より現実的な濃度分布を再現するために
% プルームモデルをシミュレーションに導入し, 
% 距離方向に連続的に変化する濃度分布に対して
% 距離分解能の要求値を検証した. 
% 距離分解能の限界値である\qty{1.5}{m}でシミュレーションしたところ, 
% 火口直上の局所的な高濃度分布に対してピークの取りこぼしが発生したが, 
% シミュレーション上の点煙源ではなく現実的な面積を持つ火口では十分に濃度分布に追従可能であると評価した. 
% また, 運用方法について火口を目視で確認し, 
% 視線と直行しないようにラマンDIALを設置することでも対応可能である. 


% ラマンDIALの受信部に用いる波長選択のための
% 干渉フィルターについて設計を行い, 
% レーザスペクトルや実行吸収断面積について検証を行った. 
% ピーク効率\qty{81}{\%}として, 
% onチャンネル側は半値幅\qty{9.2}{nm}, 
% offチャンネル側は半値幅\qty{10.6}{nm}のフィルター関数でシミュレーションし,  
% レーザ波長を十分に遮断しつつ,
% 目的とする\ce{O2}ラマン散乱信号を選択的に受信可能と評価した. 

% }
